\chapter{108}\label{section-108}

\section{Bettmeralp/VS (CH), im
Oktober}\label{bettmeralpvs-ch-im-oktober}

Louis

Bereits aus Distanz erkenne ich die längliche Öffnung auf Bodenhöhe des
Zauns.

Flexa geht hechelnd neben mir her. Zwischendurch schaut sie zu mir hoch.
Ich kann mich an ihrem Hundekopf nicht sattgucken. Auf ihrer Schnauze
liegt wieder dieses ansteckende Lächeln, das sie unermüdlich neu
anzumachen scheint. Ich streichle ihr während des Abstiegs über den
Kopf, den Ohren entlang und über ihren Rücken. Beeindruckend, wie sie
sich gestern vor \emph{«Perla»} gestellt hat. Wie standhaft sie ihren
Job macht.

Es ist eindeutig. Ich stelle die Zaunrolle neben mir ab und knie mich
nieder. Da war kräftig gebuddelt worden. Joh hatte recht. Er wird bald
hier sein. Aus Erfahrung wird er wissen, ob es tatsächlich ein Wolf war.

Ich lasse meinen Blick über das Tal gleiten. Noch immer kann ich nicht
glauben, dass mein Versteckspiel tatsächlich zu Ende ist. Jedes Mal,
wenn ich meinen wahren Namen höre, werde ich daran erinnert. Direkter
noch, wenn mich jemand aus Versehen Ciro nennt.

Ich sehe Joh kommen. Beladen mit Werkzeugen nähert er sich uns.

«Schau, Flexa, er kommt.»

Die Hündin pfeilt mit einer Leichtigkeit hoch, die ich bewundere. Wo
nimmt sie nur diese Power her? Ein deutlicher Unterschied zur
Fortbewegung der Schafe. Flexa ist eben doch ein Hund. Ich denke an Johs
Worte: «\emph{Herdenschutzhunde wachsen mit den Schafen gemeinsam auf.
Man sagt sogar, dass sie sich selbst für ein Schaf halten».} Ich
schmunzle und fühle mich gerade auch ganz und gar zugehörig.

«Hi, Louis, legen wir los. Wir müssen uns beeilen. Die Schafe werden im
Stall ungemütlich, sie drängen auf die Weide. David kümmert sich um sie.
Sobald wir hier fertig sind, wird er die Tiere zusammen mit Manuela auf
die Weide rauslassen.»

«Klar,» ich zeige auf das Loch im Zaun, «schau mal, da nützte auch der
Strom nichts. Das Tier hat offensichtlich einen schmalen Spalt entdeckt
und clever unten durch gegraben...»

Joh nickt. Er sieht mürrisch aus.

«Hhm, ja, sieht so aus,» und während er beginnt, das zusammengedrückte
Stück Zaun mit einer Eisenschere abzuzwacken, blickt er plötzlich
besorgt zu mir hinüber.

«Ich hoffe, Manuela schafft es, \emph{«Perla»} wieder
zusammenzuflicken.»

Er legt das Eisengeflecht zur Seite und stützt sich auf seinen
Oberschenkeln ab.

«Wie lange ihr Lamm Gastrecht an den Zitzen ihrer Amme hat, ist schwer
zu sagen.»

Er zwinkert mir müde zu und sieht offenbar meinen fragenden Blick.

«Fremde Lämmer können von Schafen von einem Moment auf den nächsten
einfach abgestossen werden. Wenn irgend möglich, wünschen wir unseren
Tieren, dass sie natürlich aufwachsen,» er steht auf, «dazu kommt, dass
die Flaschenfütterung um einiges aufwendiger ist.»

Während wir hoch gehen, schaue ich Joh von der Seite an. Er überrascht
mich nach wie vor fast täglich. Wie zügig das wieder ging. Seine
Handgriffe sitzen, als hätte er sich jeden einzelnen minutiös
antrainiert. Nun, hat er über die Jahre wohl auch.

«Joh, dein handwerkliches Geschick ist einfach Klasse.»

Er schielt zu mir hinüber, und zum ersten Mal glaube ich eine leichte
Verlegenheit auf seinem Gesicht zu erkennen. Vielleicht täusche ich
mich?

Wenige Schritte bevor wir oben sind, vibriert mein Handy. Valentina. Der
ersehnte Anruf. Joh gibt mir ein Zeichen.

«Nimm dir Zeit, Louis,» ruft er mir zu. Er weiss, worum es geht. Ich
setze mich mitten auf die Weide.

«Valentina?»

«Hi, Louis, schön, dass ich dich gleich erwische.»

«Ja. Wie geht es dir?»

«Hm, eigentlich ein bisschen besser. Die Dinge beginnen sich zu klären.
Ich komme langsam aus dem Tunnel. Die wochenlange Ungewissheit um dich
hat mir mehr zugesetzt, als ich dachte, Louis. Aber das weisst du.»

«Ja, Valentina, vielleicht habe ich es dir schon zu oft gesagt, aber,
glaube mir, es tut mir unglaublich leid, ich..,» und ich schlucke ein
gutes Stück Scham hinunter.

«Na ja, ich verstehe dich ja unterdessen. Darum rufe ich dich nicht an.»

«Ich weiss. Seit gestern kann ich es kaum erwarten, dass du dich
meldest. Du schreibst von erstaunlichen News?»

Mir ist, als würden die Geräte Valentinas Aufregung übertragen.

«Ja, das kann man wohl sagen. Ich war bei Giorgia im \emph{Ospedale}. Da
waren gleichzeitig auch ihre Eltern. Also, wo fang ich nur an?»

«\emph{So mach schon»} würde ich am liebsten ins Handy rufen. Nächstens
platze ich vor Anspannung.

«Nun, Giorgia ist schon fast wieder so klar, wie wir sie kennen, ausser,
ausser natürlich körperlich, sie ringt mit sich, aber eben, sie redet
und ereifert sich wie gehabt. Zusammen mit ihren Eltern hat sie mir
alles erzählt. Bastiano soll nachgewiesen werden, dass er an der Party
ziemlich scharfe Drogen gedealt hat. Gut, das tat er oft, aber, Giorgia
soll er einen heftigen Cocktail ins Getränk gemischt haben ohne ihr
Wissen. Es ist zum Kotzen. Und Giorgia erzählte, dass die Polizei sie
über einen sexuellen Übergriff durch Bastiano befragt hat. Sie hat beim
Erzählen sonderbar gelacht. Das wäre noch, meinte sie, so ein Witz.»

Ich schaffe es nicht, meine Wut zu zähmen.

«Dieses Schwein, ich könnte ihn,» ich unterbreche mich und bin froh,
dass Valentina weiterspricht.

«Du bist nicht der Einzige, Louis, wir sind alle schockiert. Bei Giorgia
wurden dieselben Substanzen im Körper gefunden wie bei Bastiano. Es soll
noch weitere Beweise geben, aber damit rückt die Polizei nicht raus. Die
wissen ja von mir, dass Giorgia mir sagte, sie glaube Basti sei bei ihr
gewesen. Sie sei sich aber nicht sicher. Gegenüber der Polizei nahm sie
die Worte nachdrücklich zurück. Sie habe keine klaren Erinnerungen. Das
Ganze ist so verworren, Louis. Und du stellst dir nicht vor, wie Bastis
Vater auf die Anschuldigungen reagiert hat.»

Und \emph{wie} ich mir das vorstellen kann.

«Doch, Valentina, kann ich. Der hat bestimmt seinen Drohhammer
ausgepackt,» ich höre mich immer lauter werden, «der will doch nur
seinen eigenen Arsch retten. Um seinen Sohn geht es ihm bestimmt nicht.»

«So ungefähr, du hast ja recht. Was ich nicht verstehe, Louis, sie
bleibt dabei, dass sie dich am \emph{Roccia} nicht gesehen hat.»

Was soll ich darauf sagen? Ich stehe da wie ein Lügner, der etwas
bezeugt, das sie verneint. Joh hat recht, es wird Zeit, dass ich mich
überwinde, Giorgia zu besuchen. Ich brauche das Gespräch mit ihr.

«Ich, ich weiss es nicht, Valentina.»

«Will sie dich vielleicht irgendwie schützen?»

Jetzt wird es gefährlich. Ich suche nach Bens Ratschlägen.

Valentina kommt mir zuvor. Zum Glück.

«Ich vermute, die Polizei geht davon aus, dass Giorgia komplett neben
sich stand und sich nicht mehr an deine Anwesenheit erinnert.»

Ich ringe mit mir. Suche schnell nach einer passenden Wendung.

«Wie schätzt du Giorgias Heilungsprozess ein?»

Es klappt.

«Sie sagt, die Ärzte seien zufrieden mit ihr. Sie mache Fortschritte.
Sie hat täglich Therapie. Es sei soweit, dass sie das \emph{Ospedale}
verlassen könne. Bis zu ihrem Austritt soll auch der Gips am Arm weg
sein. Zwischendurch sitzt sie tagsüber im Rollstuhl. Giorgia hofft,
wieder auf die Beine zu kommen. Hofft, ist gut, sie will das. Es soll
zwar möglich sein, nur wird es Monate dauern, bis sie wieder alleine
gehen kann. Wenn überhaupt. Ich,» Valentina spricht plötzlich leiser,
«ich konnte sie nicht mehr ansehen, als sie davon sprach, Louis.»

Jetzt scheint Valentina von einer Männerstimme unterbrochen zu werden.

«Ich muss Schluss machen, Louis, die nächste Vorlesung beginnt, melde
mich später wieder, ciao.»

Ich lasse mich nach hinten ins Gras fallen und strecke Arme und Beine
von mir. Valentinas plötzliches Beenden des Anrufs kam zum richtigen
Zeitpunkt. Zwischenfälle können auch rettend sein.

Ich denke an Bastiano, seine einflussreiche Familie, an ihre Macht, und
bin sicher, Bastiano wird ohne grosses Aufhebens freigekauft. Die Presse
wird bestochen oder bedroht, und weiterhin wagt es niemand, sich mit
denen anzulegen. Geht es um Ehre, um Angst oder um pure Feigheit? Was
Joh wohl darüber denkt? Ich werde ihn fragen.

Von oben erreicht mich plötzlich mehrstimmiges Blöken. Zwischendurch
Flexas Gebell. Beim Hochschauen sehe ich eine erste Gruppe Schafe
ungestüm auf die Wiese traben und sich mir nähern. Wie leicht es sein
kann, glücklich zu sein. Ich möchte ein Schaf sein.

Louis

Durchkämpfen und Dranbleiben. Joh hat ja recht.

Die Gespräche mit ihm sind ein Segen, nur, ich bin immer wieder drauf
und dran, die Termine zu verschieben oder - im Verzweiflungsmodus - ganz
abzusagen. Wenn ich an Grenzen stosse, die mir unüberwindbar erscheinen.
Wenn mich die Traurigkeit zu ersticken droht, oder wie gerade jetzt.
Wenn ich an Giorgia denke und an Johs letzte Worte gestern, bevor ich
mich zerknirscht aus dem Gespräch zurückzog. Meine häufigen Downs, die
sich als Widerstände zeigen, wie Joh meint, scheinen ihn wenig zu
beeindrucken. Manchmal ärgere ich mich über seine Ruhe, die ich in
diesen Situationen als Provokation erlebe.

Seine letzte Frage hallt in mir nach.

«Was würdest du Giorgia denn gerne sagen wollen?»

«\emph{Was wohl?»,} hätte ich ihm um ein Haar an den Kopf geworfen,
konnte mich aber gerade noch beherrschen. Unterdessen weiss ich zwar, wo
starke Gefühle aufbrechen, liegt was im Argen. Schön, aber was tue ich
damit? Was hat dieses Herumwühlen in mir am Ende für einen Nutzen? Und
der Gedanke macht mich noch wütender.

Ich ziehe meine geballten Fäuste aus den Hosentaschen, schüttle genervt
meinen Kopf und hoffe auf einen Einfall. Dann klingelt mein Handy. Die
Betriebsamkeit auf der «\emph{Sunnewaid»} hat gerade ihr Gutes. Manuela
braucht mich. Als Assistenten, meint sie, beim Wechseln des Verbandes um
\emph{Perlas} Bauch.

Während ich hochtrotte, fühle ich mich wie ein quengelndes Kind.
Unterdessen sind mir meine Stimmungsschwankungen wenigstens vertraut.

Auf dem Vorplatz der «Sunnewaid» angekommen, richte ich mich gerade auf.
«Du willst ein Mann sein», rede ich vor mich hin, «und verhältst dich
wie ein unreifer Teenager. Reiss dich zusammen.»

Vor der Praxistür ziehe ich die Schuhe aus. Nach kurzem Klopfen trete
ich ein. Als Erstes sehe ich in Manuelas strahlendem Gesicht, dass sie
sich über mein Kommen freut. Ich mag ihre ausgeglichene Art. Ganz
besonders, wenn ich mich aus meinen Tiefs hochrappeln muss.

Manuela kippt \emph{Perla} mit einem Ruck in Seitenlage. Die erschrickt,
strampelt, blökt und versucht, vom Tisch zu springen. Ich muss all meine
Kraft mobilisieren, sie festzuhalten.

«Gleich haben wir es, meine Liebe, ich will dir doch nur helfen.»

Manuela redet leise, monoton und überaus liebevoll auf das Schaf ein.
Wie fürsorglich und ruhig sie vorgeht.

Wie ich in die Augen des Tieres blicke, fallen mir plötzlich
verschiedene Situationen der letzten Wochen ein. Beim Gespräch mit Joh
zwecks Umgestaltung meines Arbeitsvertrags bleibe ich stecken. Dabei
höre ich noch einmal Marcos Worte. Endlich habe ich es geschafft, mich
bei ihm zu melden. Zu Johs Angeboten hat er mir viele Fragen gestellt
und plötzlich sein Sprechtempo gedrosselt.

«Was hast du für ein Glück, Louis, an diesen grosszügigen Menschen
geraten zu sein,» hatte er ergriffen gesagt, «und dass er sogar bereit
ist, dich persönlich und einfach obendrauf, zu begleiten. Das ist, also,
das, solche Menschen sind dünn gesät, Louis.»

Es ist, als würde gerade ein Licht angemacht. Jetzt erst glaube ich zu
verstehen, was mir Marco sagen wollte. Ein heftiges Gefühl der
Beschämung kommt hoch. Und in diesem Moment fasse ich einen Entschluss.
